% Generated mechanically from frozen input p12/ko/stacks-perfect.tex.
% Do not edit this generated file; edit the bound locale source instead.

% OK, start here.
%

\setcounter{chapter}{103}
\chapter{스택의 유도 범주}


\phantomsection
\label{stacks-perfect-section-phantom}





\section{도입}
\label{stacks-perfect-section-introduction}

\noindent
이 장에서는 대수적 스택에 딸린 유도 범주를 다룬다.
특히 이는 준연접 층의 유도 범주를 뜻한다. 즉, 스킴에 관한 결과
(스킴의 유도 범주, 절 \ref{perfect-section-introduction} 참조)와
대수공간에 관한 결과(공간의 유도 범주, 절
\ref{spaces-perfect-section-introduction} 참조)에 대응하는 결과를 증명한다.
이 장의 결과가 \cite{LM-B}의 결과와 다른 주된 이유는 우리가 일관되게
``큰 사이트''를 사용하기 때문이다. 이 장을 읽기 전에 여기서 사용하는
용어를 도입한 ``대수적 스택 위의 층'' 장과
``대수적 스택의 코호몰로지'' 장을 간단히 살펴보기 바란다.



\section{규약, 표기법 및 언어의 남용}
\label{stacks-perfect-section-conventions}

\noindent
스택의 성질, 절 \ref{stacks-properties-section-conventions}에서
도입한 규약과 언어의 남용을 계속 사용한다.
표기법은 스택의 코호몰로지, 절
\ref{stacks-cohomology-section-notation}에서 설명한 대로 사용한다.












\section{lisse-\'etale 및 flat-fppf 사이트}
\label{stacks-perfect-section-lisse-etale}

\noindent
이 절은 유도 범주에 대해 스택의 코호몰로지, 절
\ref{stacks-cohomology-section-lisse-etale}에 대응하는 내용이다.

\begin{lemma}
\label{stacks-perfect-lemma-shriek-derived}
$\mathcal{X}$를 대수적 스택이라 하자.
표기는 스택의 코호몰로지의 보조정리
\ref{stacks-cohomology-lemma-lisse-etale}와
\ref{stacks-cohomology-lemma-lisse-etale-modules}에서와 같다.
\begin{enumerate}
\item 함자
$g_! : \textit{Ab}(\mathcal{X}_{lisse,\etale}) \to
\textit{Ab}(\mathcal{X}_\etale)$에는 좌유도 함자
$$
Lg_! :
D(\mathcal{X}_{lisse,\etale})
\longrightarrow
D(\mathcal{X}_\etale)
$$
가 존재한다. 이는 $g^{-1}$의 왼쪽 수반이며 $g^{-1}Lg_! = \text{id}$를 만족한다.
\item 함자 $g_! : 
\textit{Mod}(\mathcal{X}_{lisse,\etale},
\mathcal{O}_{\mathcal{X}_{lisse,\etale}}) \to
\textit{Mod}(\mathcal{X}_\etale, \mathcal{O}_{\mathcal{X}})$에는 좌유도 함자
$$
Lg_! :
D(\mathcal{O}_{\mathcal{X}_{lisse,\etale}})
\longrightarrow
D(\mathcal{X}_\etale, \mathcal{O}_\mathcal{X})
$$
가 존재한다. 이는 $g^*$의 왼쪽 수반이며 $g^*Lg_! = \text{id}$를 만족한다.
\item 함자 $g_! : \textit{Ab}(\mathcal{X}_{flat,fppf}) \to
\textit{Ab}(\mathcal{X}_{fppf})$에는 좌유도 함자
$$
Lg_! :
D(\mathcal{X}_{flat, fppf})
\longrightarrow
D(\mathcal{X}_{fppf})
$$
가 존재한다. 이는 $g^{-1}$의 왼쪽 수반이며 $g^{-1}Lg_! = \text{id}$를 만족한다.
\item 함자 $g_! :
\textit{Mod}(\mathcal{X}_{flat,fppf},
\mathcal{O}_{\mathcal{X}_{flat,fppf}}) \to
\textit{Mod}(\mathcal{X}_{fppf}, \mathcal{O}_{\mathcal{X}})$에는 좌유도 함자
$$
Lg_! :
D(\mathcal{O}_{\mathcal{X}_{flat, fppf}})
\longrightarrow
D(\mathcal{O}_\mathcal{X})
$$
가 존재한다. 이는 $g^*$의 왼쪽 수반이며 $g^*Lg_! = \text{id}$를 만족한다.
\end{enumerate}
주의: 가군에 대한 $Lg_!$가 아벨 층에 대한 $Lg_!$와 일치하는지는
(선험적으로) 분명하지 않다. 사이트 위의 코호몰로지, 비고
\ref{sites-cohomology-remark-when-derived-shriek-equal}를 보라.
\end{lemma}

\begin{proof}
함자 $Lg_!$의 존재성과 $g^*$와의 수반성은 사이트 위의 코호몰로지,
보조정리 \ref{sites-cohomology-lemma-existence-derived-lower-shriek}이다.
(아벨 층의 경우에는 상수층 $\mathbf{Z}$를 구조층으로 사용한다.)
또한 복합체 $\mathcal{H}^\bullet$에 대한 그 값은 적절한 왼쪽 분해
$\mathcal{K}^\bullet \to \mathcal{H}^\bullet$를 취하고
$\mathcal{K}^\bullet$에 함자 $g_!$를 적용하여 계산한다.
스택의 코호몰로지의 보조정리
\ref{stacks-cohomology-lemma-lisse-etale-modules}와
\ref{stacks-cohomology-lemma-lisse-etale}에 의해
$g^{-1}g_!\mathcal{K}^\bullet = \mathcal{K}^\bullet$이므로,
각 경우에 마지막 주장이 성립함을 알 수 있다.
\end{proof}

\begin{lemma}
\label{stacks-perfect-lemma-lisse-etale-functorial-derived}
가정과 표기는 스택의 코호몰로지, 보조정리
\ref{stacks-cohomology-lemma-lisse-etale-functorial}에서와 같다고 하자.
비유계 유도 범주 위에서
$$
g^{-1} \circ Rf_* = Rf'_* \circ (g')^{-1}
\quad\text{및}\quad
L(g')_! \circ (f')^{-1} = f^{-1} \circ Lg_!
$$
가 성립한다(가군의 경우와 아벨 층의 경우 모두).
\end{lemma}

\begin{proof}
$\tau = \etale$라 하자(각각\ $\tau = fppf$라 하자). $\mathcal{F}$를
$\mathcal{X}_\tau$ 위의 아벨 층이라 하자. 스택의 코호몰로지, 보조정리
\ref{stacks-cohomology-lemma-lisse-etale-functorial-cohomology}에 의해
표준적인 (밑변환) 사상
$$
g^{-1}Rf_*\mathcal{F} \longrightarrow Rf'_*(g')^{-1}\mathcal{F}
$$
은 동형이다. 나머지 증명은 형식적이다. 아벨 군의 코호몰로지와
가군층의 코호몰로지가 일치하므로, $\mathcal{F}$가
$\mathcal{X}_\tau$ 위의 가군층인 경우에도
$g^{-1} Rf_*\mathcal{F} = Rf'_* (g')^{-1}\mathcal{F}$임을 얻는다.

\medskip\noindent
다음으로 $\mathcal{Y}_{lisse,\etale}$ 위의 (가군층 또는
아벨 군층) $\mathcal{G}$에 대해(각각\ 
$\mathcal{Y}_{flat,fppf}$ 위에서) 표준 사상
$$
L(g')_!(f')^{-1}\mathcal{G} \to f^{-1}Lg_!\mathcal{G}
$$
이 동형임을 보이자. 이를 위해서는 $\mathcal{X}_\tau$ 위의 임의의
단사층 $\mathcal{I}$에 대해 유도된 사상
$$
\Hom(L(g')_!(f')^{-1}\mathcal{G}, \mathcal{I}[n])
\leftarrow
\Hom(f^{-1}Lg_!\mathcal{G}, \mathcal{I}[n])
$$
이 모든 $n \in \mathbf{Z}$에 대해 동형임을 보이면 충분하다. (Hom은
적절한 유도 범주에서 취한다.) $f^{-1}$과 $Rf_*$의 수반성,
$Lg_!$와 $g^{-1}$의 수반성 및 이들의 ``프라임 붙은'' 판본에 의해,
이는 위에서 증명한 동형
$g^{-1} Rf_*\mathcal{I} \to Rf'_* (g')^{-1}\mathcal{I}$에서 따른다.

\medskip\noindent
$\mathcal{Y}_{lisse,\etale}$ 위의 유계 복합체 $\mathcal{G}^\bullet$
(가군 또는 아벨 군의 복합체)의 경우(각각\ 
$\mathcal{Y}_{fppf}$ 위에서), 표준 사상
\begin{equation}
\label{stacks-perfect-equation-to-show}
L(g')_!(f')^{-1}\mathcal{G}^\bullet \to f^{-1}Lg_!\mathcal{G}^\bullet
\end{equation}
은 동형이다. 이는 절단을 사용하는 통상적인 논증과
$L(g')_!(f')^{-1}$ 및 $f^{-1}Lg_!$가 삼각 범주의 완전 함자라는
사실을 적용하여 층의 경우로부터 따른다.

\medskip\noindent
$\mathcal{G}^\bullet$이 $\mathcal{Y}_{lisse,\etale}$ 위의 위로 유계인
복합체(가군 또는 아벨 군의 복합체)라고 하자(각각\ 
$\mathcal{Y}_{fppf}$ 위에서). 표준 사상 (\ref{stacks-perfect-equation-to-show})은
동형이다. 실제로 우둔한 절단 $\sigma_{\geq -n}$
(호몰로지, 절 \ref{homology-section-truncations} 참조)을 사용하여
$\mathcal{G}^\bullet$을 유계 복합체들의 쌍대극한
$\mathcal{G}^\bullet = \colim \mathcal{G}_n^\bullet$으로 쓸 수 있다.
이로부터 특수 삼각형
$$
\bigoplus\nolimits_{n \geq 1} \mathcal{G}_n^\bullet \to
\bigoplus\nolimits_{n \geq 1} \mathcal{G}_n^\bullet \to
\mathcal{G}^\bullet \to \ldots
$$
을 얻으며, 함자 $L(g')_!$, $(f')^{-1}$, $f^{-1}$, $Lg_!$ 각각은
(복합체의) 직합과 가환한다.

\medskip\noindent
$\mathcal{G}^\bullet$이 $\mathcal{Y}_{lisse,\etale}$ 위의 임의의
복합체(가군 또는 아벨 군의 복합체)라면(각각\ 
$\mathcal{Y}_{fppf}$ 위에서), 표준 절단 $\tau_{\leq n}$
(호몰로지, 절 \ref{homology-section-truncations} 참조)을 사용하여
$\mathcal{G}^\bullet$을 위로 유계인 복합체들의 쌍대극한으로 쓰고
앞 문단의 논증을 반복한다.

\medskip\noindent
마지막으로 $f^{-1}$과 $Rf_*$의 수반성, $Lg_!$와 $g^{-1}$의 수반성,
그리고 이들의 ``프라임 붙은'' 판본에 의해, 보조정리의 첫 번째 항등식은
완전한 일반성 아래 두 번째 항등식으로부터 따른다.
\end{proof}

\begin{lemma}
\label{stacks-perfect-lemma-higher-shriek-quasi-coherent}
$\mathcal{X}$를 대수적 스택이라 하자. 표기는 스택의 코호몰로지,
보조정리 \ref{stacks-cohomology-lemma-lisse-etale}에서와 같다.
\begin{enumerate}
\item $\mathcal{H}$를 $\mathcal{X}$의 lisse-\'etale 사이트 위의 준연접
$\mathcal{O}_{\mathcal{X}_{lisse,\etale}}$-가군이라 하자.
모든 $p \in \mathbf{Z}$에 대해 층 $H^p(Lg_!\mathcal{H})$는
$\mathcal{X}$ 위에서 평탄 밑변환 성질을 갖는 국소 준연접 가군이다.
\item $\mathcal{H}$를 $\mathcal{X}$의 flat-fppf 사이트 위의 준연접
$\mathcal{O}_{\mathcal{X}_{flat,fppf}}$-가군이라 하자.
모든 $p \in \mathbf{Z}$에 대해 층 $H^p(Lg_!\mathcal{H})$는
$\mathcal{X}$ 위에서 평탄 밑변환 성질을 갖는 국소 준연접 가군이다.
\end{enumerate}
\end{lemma}

\begin{proof}
스킴 $U$와 전사 매끄러운 사상 $x : U \to \mathcal{X}$를 택하자.
사이트 위의 가군, 정의 \ref{sites-modules-definition-site-local}에 의해,
각 당김 $f_i^{-1}\mathcal{H}$가 전역 표시를 갖도록 하는
\'etale 피복(각각\ fppf 피복) $\{U_i \to U\}_{i \in I}$가 존재한다
(사이트 위의 가군, 정의 \ref{sites-modules-definition-global} 참조).
여기서 $f_i : U_i \to \mathcal{X}$는 대수적 스택의 사상인 합성
$U_i \to U \to \mathcal{X}$이다.
(당김은 $\mathcal{X}/f_i$로의 제한과 ``같음''을 상기하라. 스택 위의 층,
정의 \ref{stacks-sheaves-definition-pullback} 및 그 뒤의 논의를 보라.)
피복을 세분하여 각 $U_i$가 아핀 스킴이라고 가정할 수 있다.
각 $f_i$가 매끄럽기(각각\ 평탄하기) 때문에
보조정리 \ref{stacks-perfect-lemma-lisse-etale-functorial-derived}에 의해
$f_i^{-1}Lg_!\mathcal{H} = Lg_{i, !}(f'_i)^{-1}\mathcal{H}$이다.
스택의 코호몰로지, 보조정리
\ref{stacks-cohomology-lemma-check-lqc-fbc-on-covering}을 사용하면,
보조정리의 주장은 $\mathcal{H}$가 전역 표시를 갖고 어떤 아핀 스킴
$X = \Spec(A)$에 대해 $\mathcal{X} = (\Sch/X)_{fppf}$인 경우로 귀착된다.

\medskip\noindent
우리의 표시가 다음과 같다고 하자.
$$
\bigoplus\nolimits_{j \in J} \mathcal{O} \longrightarrow
\bigoplus\nolimits_{i \in I} \mathcal{O} \longrightarrow
\mathcal{H} \longrightarrow 0
$$
여기서 $\mathcal{O} = \mathcal{O}_{\mathcal{X}_{lisse,\etale}}$이고
(각각\ $\mathcal{O} = \mathcal{O}_{\mathcal{X}_{flat,fppf}}$이다).
사이트 $\mathcal{X}_{lisse,\etale}$
(각각\ $\mathcal{X}_{flat,fppf}$)에는 끝 대상, 즉 준콤팩트한
$X/X$가 있음을 주목하라(사이트 위의 코호몰로지, 절
\ref{sites-cohomology-section-limits} 참조). 따라서
$$
\Gamma(\bigoplus\nolimits_{i \in I} \mathcal{O}) =
\bigoplus\nolimits_{i \in I} A
$$
가 사이트, 보조정리 \ref{sites-lemma-directed-colimits-sections}에 의해
성립한다. 따라서 표시의 사상은 $A$-가군 $M$의 유사한 표시
$$
\bigoplus\nolimits_{j \in J} A \longrightarrow
\bigoplus\nolimits_{i \in I} A \longrightarrow
M \longrightarrow 0
$$
에 대응한다. 또한 $\mathcal{H}$는 $M$에 딸린 준연접 층 $M^a$를
lisse-\'etale(각각\ flat-fppf) 사이트에 제한한 것과 같다. 자유
$A$-가군에 의한 분해
$$
\ldots \to F_2 \to F_1 \to F_0 \to M \to 0
$$
를 택하자. 복합체
$$
\ldots \mathcal{O} \otimes_A F_2 \to \mathcal{O} \otimes_A F_1 \to
\mathcal{O} \otimes_A F_0 \to \mathcal{H} \to 0
$$
는 자유 $\mathcal{O}$-가군에 의한 $\mathcal{H}$의 분해이다. 실제로
$\mathcal{X}_{lisse,\etale}$의 각 대상 $U/X$에 대해
(각각\ $\mathcal{X}_{flat,fppf}$에서) 구조 사상 $U \to X$가 평탄하다.
따라서 구성에 의해 $Lg_!\mathcal{H}$의 값은
$$
\ldots \to
\mathcal{O}_\mathcal{X} \otimes_A F_2 \to
\mathcal{O}_\mathcal{X} \otimes_A F_1 \to
\mathcal{O}_\mathcal{X} \otimes_A F_0 \to 0 \to \ldots
$$
이다. 이는 $\mathcal{X}_\etale$ 위의 준연접 가군의 복합체이므로
(각각\ $\mathcal{X}_{fppf}$ 위에서), 스택의 코호몰로지, 명제
\ref{stacks-cohomology-proposition-loc-qcoh-flat-base-change}에 의해
$H^p(Lg_!\mathcal{H})$는 준연접이다.
\end{proof}




\section{코호몰로지와 lisse-\'etale 및 flat-fppf 사이트}
\label{stacks-perfect-section-compare-unbounded-lisse-etale}

\noindent
대수적 스택 $\mathcal{X}$ 위의 층의 코호몰로지를 flat-fppf 사이트에서
계산할 수 있음을 이미 보았다. 이 절에서는 $\mathcal{X}$의 직접 범주의
(아마도) 비유계 대상에 대해서도 같은 사실이 참임을 증명한다. % P12-ERR-0041: frozen source wording preserved.

\begin{lemma}
\label{stacks-perfect-lemma-higher-shriek-Z}
$\mathcal{X}$를 대수적 스택이라 하자. 보조정리
\ref{stacks-perfect-lemma-shriek-derived}의 (1)에 나오는 $Lg_!$ 또는
보조정리 \ref{stacks-perfect-lemma-shriek-derived}의 (3)에 나오는 $Lg_!$ 어느 경우에도
$Lg_!\mathbf{Z} = \mathbf{Z}$이다.
\end{lemma}

\begin{proof}
flat-fppf 사이트와 fppf 사이트의 비교에 대해 이를 증명하겠다.
lisse-\'etale 사이트의 경우도 완전히 같다.
$i \not = 0$일 때 $H^i(Lg_!\mathbf{Z})$가 $0$이고 표준 사상
$H^0(Lg_!\mathbf{Z}) \to \mathbf{Z}$가 동형임을 보여야 한다.
$\mathcal{U}$가 스킴이고 $f$가 또한 국소 유한 표시인 전사 평탄 사상
$f : \mathcal{U} \to \mathcal{X}$를 택하자.
(예를 들어 표시 $U \to \mathcal{X}$를 택하고 $\mathcal{U}$를
$U$에 대응하는 대수적 스택으로 둔다.)
스택 위의 층, 보조정리
\ref{stacks-sheaves-lemma-check-exactness-covering}와
\ref{stacks-sheaves-lemma-surjective-flat-locally-finite-presentation}에 의해,
$i \not = 0$일 때 당김 $f^{-1}H^i(Lg_!\mathbf{Z})$가 $0$이고
$H^0(Lg_!\mathbf{Z}) \to f^{-1}\mathbf{Z}$의 당김이 동형임을
보이는 것으로 충분하다. 보조정리
\ref{stacks-perfect-lemma-lisse-etale-functorial-derived}에 의해
$f^{-1}Lg_!\mathbf{Z} = L(g')_!\mathbf{Z}$를 얻는다. 여기서
$g' : \Sh(\mathcal{U}_{flat, fppf}) \to \Sh(\mathcal{U}_{fppf})$
는 $\mathcal{U}$에 대한 해당 비교 사상이다.
이로써 다음 문단에서 다루는 경우로 귀착된다.

\medskip\noindent
어떤 스킴 $X$에 대해 $\mathcal{X} = (\Sch/X)_{fppf}$라고 하자.
이 경우 범주 $\mathcal{X}_{flat, fppf}$에는 끝 대상 $e$, 즉 $X/X$가 있고,
더욱이 함자 $u : \mathcal{X}_{flat, fppf} \to \mathcal{X}_{fppf}$는
$e$를 끝 대상으로 보낸다. $\mathbf{Z}$는 끝 대상 위의 자유 아벨 층이므로
(끝 대상이 존재한다면), 사이트 위의 코호몰로지, 보조정리
\ref{sites-cohomology-lemma-existence-derived-lower-shriek}에 나오는
$Lg_!$의 바로 그 구성에 의해 $Lg_!\mathbf{Z} = \mathbf{Z}$를 얻는다.
\end{proof}

\begin{lemma}
\label{stacks-perfect-lemma-lisse-etale-cohomology}
$\mathcal{X}$를 대수적 스택이라 하자. 표기는
보조정리 \ref{stacks-perfect-lemma-shriek-derived}에서와 같다.
\begin{enumerate}
\item $D(\mathcal{X}_\etale)$의 $K$에 대해 다음이 성립한다.
\begin{enumerate}
\item $R\Gamma(\mathcal{X}_\etale, K) =
R\Gamma(\mathcal{X}_{lisse,\etale}, g^{-1}K)$이고,
\item $\mathcal{X}_{lisse,\etale}$의 임의의 대상 $x$에 대해
$R\Gamma(x, K) =
R\Gamma(\mathcal{X}_{lisse,\etale}/x, g^{-1}K)$이다.
\end{enumerate}
\item $D(\mathcal{X}_{fppf})$의 $K$에 대해 다음이 성립한다.
\begin{enumerate}
\item $R\Gamma(\mathcal{X}_{fppf}, K) =
R\Gamma(\mathcal{X}_{flat,fppf}, g^{-1}K)$이고,
\item $\mathcal{X}_{flat,fppf}$의 임의의 대상 $x$에 대해
$H^p(x, K) =
R\Gamma(\mathcal{X}_{flat,fppf}/x, g^{-1}K)$이다.
\end{enumerate}
\end{enumerate}
두 경우 모두 가군에 대해서도 같은 사실이 성립한다. 실제로
$g^{-1} = g^*$이고, 사이트 위의 코호몰로지, 보조정리
\ref{sites-cohomology-lemma-modules-abelian-unbounded}에 의해
코호몰로지를 계산할 때 차이가 없다.
\end{lemma}

\begin{proof}
flat-fppf 사이트와 fppf 사이트의 비교에 대해 이를 증명하겠다.
lisse-\'etale 사이트의 경우도 완전히 같다. 보조정리
\ref{stacks-perfect-lemma-higher-shriek-Z}에 의해 $Lg_!\mathbf{Z} = \mathbf{Z}$이다.
따라서
\begin{align*}
R\Gamma(\mathcal{X}_{fppf}, K)
& =
R\Hom(\mathbf{Z}, K) \\
& =
R\Hom(Lg_!\mathbf{Z}, K) \\
& =
R\Hom(\mathbf{Z}, g^{-1}K) \\
& =
R\Gamma(\mathcal{X}_{lisse,\etale}, g^{-1}K)
\end{align*}
를 얻는다. 이것으로 (1)(a)를 증명했다.
(1)(b)는 (1)(a)에서 따른다. 실제로 $x$가 스킴 $U$ 위에 놓이면,
사이트 $\mathcal{X}_\etale/x$는 $(\Sch/U)_\etale$와 동치이고
$\mathcal{X}_{lisse,\etale}$는 $U_{lisse, \etale}$와 동치이다.
\end{proof}









\section{준연접 가군의 유도 범주}
\label{stacks-perfect-section-derived}

\noindent
$\mathcal{X}$를 대수적 스택이라 하자. 포함 함자
$\QCoh(\mathcal{O}_\mathcal{X}) \to
\textit{Mod}(\mathcal{O}_\mathcal{X})$가 완전하지 않으므로,
$D_\QCoh(\mathcal{O}_\mathcal{X})$를 준연접 코호몰로지 층을 갖는
$D(\mathcal{O}_\mathcal{X})$의 복합체들로 이루어진 충만한 부분범주로
정의할 수 없다. 대신 스택의 코호몰로지, 비고
\ref{stacks-cohomology-remark-bousfield-colocalization}와 유사하게
준연접 가군의 유도 범주를 몫으로 정의한다.

\medskip\noindent
$\textit{LQCoh}^{fbc}(\mathcal{O}_\mathcal{X}) \subset
\textit{Mod}(\mathcal{O}_\mathcal{X})$가 평탄 밑변환 성질을 갖는
국소 준연접 $\mathcal{O}_\mathcal{X}$-가군들의 충만한 부분범주를
뜻함을 상기하라. 스택의 코호몰로지, 절
\ref{stacks-cohomology-section-loc-qcoh-flat-base-change}를 보라.
다음과 같이 줄여 쓰겠다.
$$
D_{\textit{LQCoh}^{fbc}}(\mathcal{O}_\mathcal{X}) =
D_{\textit{LQCoh}^{fbc}(\mathcal{O}_\mathcal{X})}(\mathcal{O}_\mathcal{X})
$$
유도 범주, 보조정리 \ref{derived-lemma-cohomology-in-serre-subcategory}와
스택의 코호몰로지, 명제
\ref{stacks-cohomology-proposition-loc-qcoh-flat-base-change}의 (2)로부터
$D_{\textit{LQCoh}^{fbc}}(\mathcal{O}_\mathcal{X})$가
$D(\mathcal{O}_\mathcal{X})$의 엄밀히 충만하고 포화된
삼각 부분범주임을 얻는다.

\medskip\noindent
$\textit{Parasitic}(\mathcal{O}_\mathcal{X}) \subset
\textit{Mod}(\mathcal{O}_\mathcal{X})$가 기생
$\mathcal{O}_\mathcal{X}$-가군들의 충만한 부분범주를 뜻한다고 하자.
스택의 코호몰로지, 절
\ref{stacks-cohomology-section-parasitic}를 보라.
다음과 같이 줄여 쓰겠다.
$$
D_{\textit{Parasitic}}(\mathcal{O}_\mathcal{X}) =
D_{\textit{Parasitic}(\mathcal{O}_\mathcal{X})}(\mathcal{O}_\mathcal{X})
$$
앞에서와 같이 이는 $D(\mathcal{O}_\mathcal{X})$의 엄밀히 충만하고
포화된 삼각 부분범주이다. 실제로
$\textit{Parasitic}(\mathcal{O}_\mathcal{X})$는
$\textit{Mod}(\mathcal{O}_\mathcal{X})$의 Serre 부분범주이다. 스택의
코호몰로지, 보조정리 \ref{stacks-cohomology-lemma-parasitic}를 보라.

\medskip\noindent
$\textit{Mod}(\mathcal{O}_\mathcal{X})$의 두 약한 Serre 부분범주의 교집합
$\textit{Parasitic}(\mathcal{O}_\mathcal{X}) \cap
\textit{LQCoh}^{fbc}(\mathcal{O}_\mathcal{X})$도
약한 Serre 부분범주이다.
마찬가지로 다음과 같이 줄여 쓰자.
\begin{align*}
D_{\textit{Parasitic} \cap \textit{LQCoh}^{fbc}}(\mathcal{O}_\mathcal{X})
& =
D_{\textit{Parasitic}(\mathcal{O}_\mathcal{X}) \cap
\textit{LQCoh}^{fbc}(\mathcal{O}_\mathcal{X})}(\mathcal{O}_\mathcal{X}) \\
& =
D_{\textit{Parasitic}}(\mathcal{O}_\mathcal{X})
\cap
D_{\textit{LQCoh}^{fbc}}(\mathcal{O}_\mathcal{X})
\end{align*}
앞에서와 같이 이는 $D(\mathcal{O}_\mathcal{X})$의 엄밀히 충만하고
포화된 삼각 부분범주이다. 따라서 특히
$D_{\textit{LQCoh}^{fbc}}(\mathcal{O}_\mathcal{X})$의 엄밀히 충만하고
포화된 삼각 부분범주이다.

\begin{definition}
\label{stacks-perfect-definition-derived}
$\mathcal{X}$를 대수적 스택이라 하자. 위의 표기 아래, {\it 준연접
코호몰로지 층을 갖는 $\mathcal{O}_\mathcal{X}$-가군의 유도 범주}를
Verdier 몫으로 정의한다.\footnote{이 정의는 문헌의 정의와 다르다.
\cite[6.3]{olsson_sheaves}를 보라. 그러나 보조정리
\ref{stacks-perfect-lemma-derived-quasi-coherent}에 의해 그 정의와 일치한다.}
$$
D_\QCoh(\mathcal{O}_\mathcal{X}) =
D_{\textit{LQCoh}^{fbc}}(\mathcal{O}_\mathcal{X})/
D_{\textit{Parasitic} \cap \textit{LQCoh}^{fbc}}(\mathcal{O}_\mathcal{X})
$$
\end{definition}

\noindent
Verdier 몫은 유도 범주, 절 \ref{derived-section-quotients}에서 정의된다.
$D_{\textit{LQCoh}^{fbc}}(\mathcal{O}_\mathcal{X})$의 사상
$a : E \to E'$가 $D_\QCoh(\mathcal{O}_\mathcal{X})$에서 동형이 되는 것은
그 원뿔 $C(a)$가 기생 코호몰로지 층을 갖는 것과 동치이다. 유도 범주,
보조정리 \ref{derived-lemma-operations}를 보라.

\medskip\noindent
함자들을 생각하자.
$$
D_{\textit{LQCoh}^{fbc}}(\mathcal{O}_\mathcal{X})
\xrightarrow{H^i}
\textit{LQCoh}^{fbc}(\mathcal{O}_\mathcal{X})
\xrightarrow{Q}
\QCoh(\mathcal{O}_\mathcal{X})
$$
$Q$가 부분범주
$\textit{Parasitic}(\mathcal{O}_\mathcal{X}) \cap
\textit{LQCoh}^{fbc}(\mathcal{O}_\mathcal{X})$를 소멸시킴을 주목하라.
스택의 코호몰로지, 보조정리
\ref{stacks-cohomology-lemma-adjoint-kernel-parasitic}를 보라.
유도 범주, 보조정리
\ref{derived-lemma-universal-property-quotient}에 의해 코호몰로지 함자
\begin{equation}
\label{stacks-perfect-equation-Hi-quasi-coherent}
H^i :
D_\QCoh(\mathcal{O}_\mathcal{X})
\longrightarrow
\QCoh(\mathcal{O}_\mathcal{X})
\end{equation}
를 얻는다. 또한 $E \in D_\QCoh(\mathcal{O}_\mathcal{X})$가 영 대상인 것은
모든 $i \in \mathbf{Z}$에 대해 $H^i(E) = 0$인 것과 동치임을 주목하라.
실제로 스택의 코호몰로지, 보조정리
\ref{stacks-cohomology-lemma-adjoint-kernel-parasitic}에 의해 $Q$의 핵은 정확히
$\textit{Parasitic}(\mathcal{O}_\mathcal{X}) \cap
\textit{LQCoh}^{fbc}(\mathcal{O}_\mathcal{X})$와 같다.

\medskip\noindent
범주 $\textit{Parasitic}(\mathcal{O}_\mathcal{X}) \cap
\textit{LQCoh}^{fbc}(\mathcal{O}_\mathcal{X})$와
$\textit{LQCoh}^{fbc}(\mathcal{O}_\mathcal{X})$는 또한
\'etale 위상에서의 가군의 아벨 범주
$\textit{Mod}(\mathcal{X}_\etale, \mathcal{O}_\mathcal{X})$의
약한 Serre 부분범주이다. 스택의 코호몰로지의 명제
\ref{stacks-cohomology-proposition-loc-qcoh-flat-base-change}와 보조정리
\ref{stacks-cohomology-lemma-parasitic}를 보라.
따라서 다음 보조정리의 명제가 의미를 갖는다.

\begin{lemma}
\label{stacks-perfect-lemma-compare-etale-fppf-QCoh}
$\mathcal{X}$를 대수적 스택이라 하자. 다음과 같이 줄여 쓰자.
$\mathcal{P}_\mathcal{X} = \textit{Parasitic}(\mathcal{O}_\mathcal{X}) \cap
\textit{LQCoh}^{fbc}(\mathcal{O}_\mathcal{X})$.
비교 사상 $\epsilon : \mathcal{X}_{fppf} \to \mathcal{X}_\etale$은
가환 그림
$$
\xymatrix{
D_{\textit{Parasitic} \cap \textit{LQCoh}^{fbc}}(\mathcal{O}_\mathcal{X})
\ar[r] &
D_{\textit{LQCoh}^{fbc}}(\mathcal{O}_\mathcal{X}) \ar[r] &
D(\mathcal{O}_\mathcal{X}) \\
D_{\mathcal{P}_\mathcal{X}}(\mathcal{X}_\etale, \mathcal{O}_\mathcal{X})
\ar[r] \ar[u]^{\epsilon^*} &
D_{\textit{LQCoh}^{fbc}(\mathcal{O}_\mathcal{X})}(
\mathcal{X}_\etale, \mathcal{O}_\mathcal{X})
\ar[r] \ar[u]^{\epsilon^*} &
D(\mathcal{X}_\etale, \mathcal{O}_\mathcal{X})
\ar[u]^{\epsilon^*}
}
$$
을 유도한다. 또한 왼쪽의 두 수직 화살표는 삼각 범주의 동치이므로,
다음 동치도 얻는다.
$$
D_{\textit{LQCoh}^{fbc}(\mathcal{O}_\mathcal{X})}
(\mathcal{X}_\etale, \mathcal{O}_\mathcal{X})
/
D_{\mathcal{P}_\mathcal{X}}(\mathcal{X}_\etale, \mathcal{O}_\mathcal{X})
\longrightarrow
D_\QCoh(\mathcal{O}_\mathcal{X})
$$
\end{lemma}

\begin{proof}
$\epsilon^*$는 완전하므로 보조정리의 명제에 나오는 그림을 얻는 것은
명백하다. 다음 상황에 사이트 위의 코호몰로지, 보조정리
\ref{sites-cohomology-lemma-compare-topologies-derived-adequate-modules}를
적용하여 가운데 수직 화살표가 동치임을 보이겠다.
$\mathcal{C} = \mathcal{X}$,
$\tau = fppf$,
$\tau' = \etale$,
$\mathcal{O} = \mathcal{O}_\mathcal{X}$,
$\mathcal{A} = \textit{LQCoh}^{fbc}(\mathcal{O}_\mathcal{X})$이고,
$\mathcal{B}$는 아핀 스킴 위에 놓이는 $\mathcal{X}$의 대상들의 집합이다.
보조정리를 적용할 수 있음을 확인하려면 조건 (1), (2), (3), (4)를
확인해야 한다. 조건 (1)과 (2)는 위의 논의에서 명백하다(구체적으로는
스택의 코호몰로지, 명제
\ref{stacks-cohomology-proposition-loc-qcoh-flat-base-change}에서 따른다).
모든 스킴은 아핀 열린집합들에 의한 Zariski 열린 피복을 가지므로 조건
(3)이 성립한다. 조건 (4)는 강하, 보조정리
\ref{descent-lemma-quasi-coherent-and-flat-base-change}에서 따른다.

\medskip\noindent
범주의 동치
$\epsilon^* : 
D_{\textit{LQCoh}^{fbc}(\mathcal{O}_\mathcal{X})}(
\mathcal{X}_\etale, \mathcal{O}_\mathcal{X})
\to
D_{\textit{LQCoh}^{fbc}}(\mathcal{O}_\mathcal{X})$
가 기생 코호몰로지 층을 갖는 복합체들의 부분범주 사이의 동치를
유도한다는 확인은 생략한다.
\end{proof}

\noindent
$\mathcal{X}$를 대수적 스택이라 하자.
스택의 코호몰로지, 보조정리
\ref{stacks-cohomology-lemma-quasi-coherent-weak-serre}에 의해,
준연접 가군의 범주
$\QCoh(\mathcal{O}_{\mathcal{X}_{lisse,\etale}})$는
$\textit{Mod}(\mathcal{O}_{\mathcal{X}_{lisse,\etale}})$의
약한 Serre 부분범주를 이루고, 준연접 가군의 범주
$\QCoh(\mathcal{O}_{\mathcal{X}_{flat,fppf}})$는
$\textit{Mod}(\mathcal{O}_{\mathcal{X}_{flat,fppf}})$의
약한 Serre 부분범주를 이룬다. 따라서 다음을 생각할 수 있다.
$$
D_\QCoh(\mathcal{O}_{\mathcal{X}_{lisse,\etale}}) =
D_{\QCoh(\mathcal{O}_{\mathcal{X}_{lisse,\etale}})}(
\mathcal{O}_{\mathcal{X}_{lisse,\etale}})
\subset
D(\mathcal{O}_{\mathcal{X}_{lisse,\etale}})
$$
그리고 마찬가지로
$$
D_\QCoh(\mathcal{O}_{\mathcal{X}_{flat,fppf}}) =
D_{\QCoh(\mathcal{O}_{\mathcal{X}_{flat,fppf}})}(
\mathcal{O}_{\mathcal{X}_{flat,fppf}})
\subset
D(\mathcal{O}_{\mathcal{X}_{flat,fppf}})
$$
앞에서와 같이 이들은 엄밀히 충만하고 포화된 삼각 부분범주이다.
$D_\QCoh(\mathcal{O}_\mathcal{X})$는 이들 각각과 동치임이 밝혀진다.

\begin{lemma}
\label{stacks-perfect-lemma-derived-quasi-coherent}
$\mathcal{X}$를 대수적 스택이라 하자. 다음과 같이 두자.
$\mathcal{P}_\mathcal{X} = \textit{Parasitic}(\mathcal{O}_\mathcal{X}) \cap
\textit{LQCoh}^{fbc}(\mathcal{O}_\mathcal{X})$.
\begin{enumerate}
\item
$\mathcal{F}^\bullet$을
$D_{\textit{LQCoh}^{fbc}(\mathcal{O}_\mathcal{X})}
(\mathcal{X}_\etale, \mathcal{O}_\mathcal{X})$의 대상이라 하자.
lisse-\'etale 사이트에 대해 스택의 코호몰로지, 보조정리
\ref{stacks-cohomology-lemma-lisse-etale}에 나오는 $g$를 택하면
다음이 성립한다.
\begin{enumerate}
\item $g^*\mathcal{F}^\bullet$은
$D_\QCoh(\mathcal{O}_{\mathcal{X}_{lisse,\etale}})$에 속한다.
\item $g^*\mathcal{F}^\bullet = 0$인 것은
$\mathcal{F}^\bullet$이
$D_{\mathcal{P}_\mathcal{X}}(\mathcal{X}_\etale, \mathcal{O}_\mathcal{X})$
에 속하는 것과 동치이다.
\item $\mathcal{H}^\bullet$이
$D_\QCoh(\mathcal{O}_{\mathcal{X}_{lisse,\etale}})$에 속하면
$Lg_!\mathcal{H}^\bullet$은
$D_{\textit{LQCoh}^{fbc}(\mathcal{O}_\mathcal{X})}(
\mathcal{X}_\etale, \mathcal{O}_\mathcal{X})$에 속한다.
\item 함자 $g^*$와 $Lg_!$는 서로 역인 함자들을 정의한다.
$$
\xymatrix{
D_\QCoh(\mathcal{O}_\mathcal{X}) \ar@<1ex>[r]^-{g^*} &
D_\QCoh(\mathcal{O}_{\mathcal{X}_{lisse,\etale}})
\ar@<1ex>[l]^-{Lg_!}
}
$$
\end{enumerate}
\item
$\mathcal{F}^\bullet$을
$D_{\textit{LQCoh}^{fbc}}(\mathcal{O}_\mathcal{X})$의 대상이라 하자.
flat-fppf 사이트에 대해 스택의 코호몰로지, 보조정리
\ref{stacks-cohomology-lemma-lisse-etale}에 나오는 $g$를 택하면
다음이 성립한다.
\begin{enumerate}
\item $g^*\mathcal{F}^\bullet$은
$D_\QCoh(\mathcal{O}_{\mathcal{X}_{flat, fppf}})$에 속한다.
\item $g^*\mathcal{F}^\bullet = 0$인 것은
$\mathcal{F}^\bullet$이
$D_{\mathcal{P}_\mathcal{X}}(\mathcal{O}_\mathcal{X})$에 속하는 것과
동치이다.
\item $\mathcal{H}^\bullet$이
$D_\QCoh(\mathcal{O}_{\mathcal{X}_{flat,fppf}})$에 속하면
$Lg_!\mathcal{H}^\bullet$은
$D_{\textit{LQCoh}^{fbc}}(\mathcal{O}_\mathcal{X})$에 속한다.
\item 함자 $g^*$와 $Lg_!$는 서로 역인 함자들을 정의한다.
$$
\xymatrix{
D_\QCoh(\mathcal{O}_\mathcal{X}) \ar@<1ex>[r]^-{g^*} &
D_\QCoh(\mathcal{O}_{\mathcal{X}_{flat,fppf}}) \ar@<1ex>[l]^-{Lg_!}
}
$$
\end{enumerate}
\end{enumerate}
\end{lemma}

\begin{proof}
함자 $g^* = g^{-1}$는 완전하므로 (1)(a), (2)(a), (1)(b), (2)(b)는
스택의 코호몰로지의 보조정리
\ref{stacks-cohomology-lemma-quasi-coherent}와
\ref{stacks-cohomology-lemma-parasitic-in-terms-flat-fppf}에서 따른다.

\medskip\noindent
(1)(c)와 (2)(c)를 증명하자.
보조정리 \ref{stacks-perfect-lemma-shriek-derived}의 $Lg_!$의 구성은(사이트 위의
코호몰로지, 보조정리
\ref{sites-cohomology-lemma-existence-derived-lower-shriek}를 통하며, 이는
다시 유도 범주, 명제 \ref{derived-proposition-left-derived-exists}를 사용한다)
$D(\mathcal{O}_{\mathcal{X}_{lisse,\etale}})$의 임의의 대상
$\mathcal{H}^\bullet$에 대한 $Lg_!$가 다음과 같이 계산됨을 보인다.
$$
Lg_!\mathcal{H}^\bullet = \colim g_!\mathcal{K}_n^\bullet =
g_! \colim \mathcal{K}_n^\bullet
$$
(항별 쌍대극한이다.) 여기서 준동형
$\colim \mathcal{K}_n^\bullet \to \mathcal{H}^\bullet$은 준동형
$\mathcal{K}_n^\bullet \to \tau_{\leq n} \mathcal{H}^\bullet$을 유도한다.
포함 함자
$$
\textit{LQCoh}^{fbc}(\mathcal{O}_\mathcal{X}) \subset
\textit{Mod}(\mathcal{X}_\etale, \mathcal{O}_\mathcal{X})
\quad\text{및}\quad
\textit{LQCoh}^{fbc}(\mathcal{O}_\mathcal{X}) \subset
\textit{Mod}(\mathcal{O}_\mathcal{X})
$$
가 여과 쌍대극한과 양립하므로, (c)를
$D_\QCoh(\mathcal{O}_{\mathcal{X}_{lisse,\etale}})$ 및
$D_\QCoh(\mathcal{O}_{\mathcal{X}_{flat,fppf}})$의 위로 유계인
복합체 $\mathcal{H}^\bullet$에 대해 증명하는 것으로 충분하다.
이 경우 $H^n(Lg_!\mathcal{H}^\bullet)$이
$\textit{LQCoh}^{fbc}(\mathcal{O}_\mathcal{X})$에 속함을 보이기 위해,
$i > m$이면 $\mathcal{H}^i = 0$이 되게 하는 정수 $m$에 관한 귀납법을
쓸 수 있다. $m < n$이면 $H^n(Lg_!\mathcal{H}^\bullet) = 0$이므로
결과가 성립한다. 일반적으로 특수 삼각형
$$
\tau_{\leq m - 1}\mathcal{H}^\bullet \to \mathcal{H}^\bullet \to
H^m(\mathcal{H}^\bullet)[-m] \to \ldots
$$
(유도 범주, 비고
\ref{derived-remark-truncation-distinguished-triangle})을 생각하고
함자 $Lg_!$를 적용한다.
$\textit{LQCoh}^{fbc}(\mathcal{O}_\mathcal{X})$는 가군 범주의 약한
Serre 부분범주이므로, 세 항 중 두 항에 대해 (c)를 증명하면 충분하다.
$Lg_!\tau_{\leq m - 1}\mathcal{H}^\bullet$에 대해서는 귀납법으로 결과를
얻고, $Lg_!H^m(\mathcal{H}^\bullet)[-m]$에 대해서는 보조정리
\ref{stacks-perfect-lemma-higher-shriek-quasi-coherent}로 결과를 얻는다.
따라서 (c)가 성립한다.

\medskip\noindent
(2)(d)를 증명하자. (2)(a)와 (2)(b)에 의해 함자
$g^{-1} = g^*$는 함자
$$
c :
D_\QCoh(\mathcal{O}_\mathcal{X})
\longrightarrow
D_\QCoh(\mathcal{O}_{\mathcal{X}_{flat, fppf}})
$$
를 유도한다. 유도 범주, 보조정리
\ref{derived-lemma-universal-property-quotient}를 보라.
따라서 다음과 같은 삼각 범주의 그림을 얻는다.
$$
\xymatrix{
D_{\textit{LQCoh}^{fbc}}(\mathcal{O}_\mathcal{X})
\ar[rd]^{g^{-1}} \ar[rr]_q & &
D_\QCoh(\mathcal{O}_\mathcal{X}) \ar[ld]^c \\
& D_\QCoh(\mathcal{O}_{\mathcal{X}_{flat, fppf}})
\ar@<1ex>[lu]^{Lg_!}
}
$$
여기서 $q$는 몫 함자이고, 안쪽 삼각형은 가환하며
$g^{-1}Lg_! = \text{id}$이다.
$D_{\textit{LQCoh}^{fbc}}(\mathcal{O}_\mathcal{X})$의 임의의 대상 $E$에
대해 사상 $a : Lg_!g^{-1}E \to E$는
$D(\mathcal{O}_{\mathcal{X}_{flat, fppf}})$에서 준동형으로 보내진다.
따라서 $a$의 원뿔은 $g^{-1}$ 아래 영으로 보내지고, (2)(b)에 의해
$q(a)$가 동형임을 알 수 있다. 그러므로 $q \circ Lg_!$는 $c$의
준역함자이다.

\medskip\noindent
lisse-\'etale 사이트의 경우에는 위와 정확히 같은 논증으로
$$
D_{\textit{LQCoh}^{fbc}(\mathcal{O}_\mathcal{X})}(
\mathcal{X}_\etale, \mathcal{O}_\mathcal{X})
/
D_{\mathcal{P}_\mathcal{X}}(
\mathcal{X}_\etale, \mathcal{O}_\mathcal{X})
$$
가 $D_\QCoh(\mathcal{O}_{\mathcal{X}_{lisse,\etale}})$와 동치임을
증명한다. 보조정리 \ref{stacks-perfect-lemma-compare-etale-fppf-QCoh}의 마지막 동치를
적용하면 증명이 끝난다.
\end{proof}

\noindent
다음 보조정리는 몫 함자
$D_{\textit{LQCoh}^{fbc}}(\mathcal{O}_\mathcal{X}) \to
D_\QCoh(\mathcal{O}_\mathcal{X})$가 왼쪽 수반을 가짐을 말해 준다.
비고 \ref{stacks-perfect-remark-QCoh-admissible}를 보라.

\begin{lemma}
\label{stacks-perfect-lemma-bousfield-colocalization}
$\mathcal{X}$를 대수적 스택이라 하자.
$E$를 $D_{\textit{LQCoh}^{fbc}}(\mathcal{O}_\mathcal{X})$의 대상이라 하자.
$D_{\textit{LQCoh}^{fbc}}(\mathcal{O}_\mathcal{X})$ 안에 표준적인
특수 삼각형
$$
E' \to E \to P \to E'[1]
$$
이 존재하여, $P$는
$D_{\textit{Parasitic} \cap \textit{LQCoh}^{fbc}}
(\mathcal{O}_\mathcal{X})$에 속하고 모든
$D_{\textit{Parasitic} \cap \textit{LQCoh}^{fbc}}(\mathcal{O}_\mathcal{X})$
의 $P'$에 대해
$$
\Hom_{D(\mathcal{O}_\mathcal{X})}(E', P') = 0
$$
이다.
\end{lemma}

\begin{proof}
스택의 코호몰로지, 절
\ref{stacks-cohomology-section-lisse-etale}에서 다룬 환 달린 토포스의 사상
$g : \Sh(\mathcal{X}_{flat, fppf}) \longrightarrow \Sh(\mathcal{X}_{fppf})$
를 생각하자. $E' = Lg_!g^*E$로 두고, $P$를 수반 사상
$E' \to E$의 원뿔로 두자. 보조정리 \ref{stacks-perfect-lemma-shriek-derived}의 (4)를
보라. 보조정리 \ref{stacks-perfect-lemma-derived-quasi-coherent}의 (2)(a)와 (2)(c)에 의해
$E'$은 $D_{\textit{LQCoh}^{fbc}}(\mathcal{O}_\mathcal{X})$에 속한다.
따라서 $P$도 $D_{\textit{LQCoh}^{fbc}}(\mathcal{O}_\mathcal{X})$에 속한다.
보조정리 \ref{stacks-perfect-lemma-shriek-derived}의 (4)에 의해
$g^*Lg_! = \text{id}$이므로 사상 $g^*E' \to g^*E$는 동형이다.
따라서 $g^*P = 0$이고, 그러므로 보조정리
\ref{stacks-perfect-lemma-derived-quasi-coherent}의 (2)(b)에 의해 $P$는
$D_{\textit{Parasitic} \cap \textit{LQCoh}^{fbc}}(\mathcal{O}_\mathcal{X})$
의 대상이다. 마지막으로
$P'$가
$D_{\textit{Parasitic} \cap \textit{LQCoh}^{fbc}}(\mathcal{O}_\mathcal{X})$
에 속할 때
$$
\Hom(E', P') = \Hom(Lg_!g^*E, P') = \Hom(g^*E, g^*P') = 0
$$
이다. 실제로 보조정리 \ref{stacks-perfect-lemma-derived-quasi-coherent}의 (2)(b)에 의해
$g^*P' = 0$이다. 특수 삼각형 $E' \to E \to P \to E'[1]$은 표준적이다
(더 정확히 말해, $E$ 위의 항등사상을 유도하는 삼각형들의 동형까지
유일하다). 유도 범주, 절 \ref{derived-section-admissible}의 논의를 보라.
\end{proof}

\begin{remark}
\label{stacks-perfect-remark-QCoh-admissible}
보조정리 \ref{stacks-perfect-lemma-bousfield-colocalization}의 결과는
$$
D_{\textit{Parasitic} \cap \textit{LQCoh}^{fbc}}(\mathcal{O}_\mathcal{X})
\subset
D_{\textit{LQCoh}^{fbc}}(\mathcal{O}_\mathcal{X})
$$
가 왼쪽 허용 가능 부분범주임을 말해 준다. 유도 범주, 절
\ref{derived-section-admissible}를 보라. 특히
$\mathcal{A} \subset D_{\textit{LQCoh}^{fbc}}(\mathcal{O}_\mathcal{X})$가
그 왼쪽 직교를 나타내면, 유도 범주, 명제
\ref{derived-proposition-summarize-admissible}에 의해 $\mathcal{A}$는
$D_{\textit{LQCoh}^{fbc}}(\mathcal{O}_\mathcal{X})$에서 오른쪽 허용
가능하고, 합성
$$
\mathcal{A} \longrightarrow
D_{\textit{LQCoh}^{fbc}}(\mathcal{O}_\mathcal{X}) \longrightarrow
D_\QCoh(\mathcal{O}_\mathcal{X})
$$
은 동치이다. 이는 $D_\QCoh(\mathcal{O}_\mathcal{X})$를
$D_{\textit{LQCoh}^{fbc}}(\mathcal{O}_\mathcal{X})$의 엄밀히 충만하고
포화된 삼각 부분범주로, 또
$D(\mathcal{X}_{fppf}, \mathcal{O}_\mathcal{X})$의 그러한 부분범주로
볼 수 있음을 뜻한다.
\end{remark}




\section{준연접 가군의 유도 직접상}
\label{stacks-perfect-section-derived-pushforward}

\noindent
위 내용의 첫 응용으로 유도 직접상을 구성한다.
예, 절 \ref{examples-section-derived-push-quasi-coherent}에서 독자는
대수적 스택의 준콤팩트 준분리 사상
$f : \mathcal{X} \to \mathcal{Y}$ 가운데 직접상 함자 $Rf_*$가 함자
$D_\QCoh(\mathcal{O}_\mathcal{X}) \to
D_\QCoh(\mathcal{O}_\mathcal{Y})$를 유도하지 않는 예를 찾을 수 있다.
따라서 아래로 유계인 복합체로 제한할 필요가 있다.

\begin{proposition}
\label{stacks-perfect-proposition-derived-direct-image-quasi-coherent}
$f : \mathcal{X} \to \mathcal{Y}$를 대수적 스택의 준콤팩트
준분리 사상이라 하자. 함자 $Rf_*$는 가환 그림
$$
\xymatrix{
D^{+}_{\textit{Parasitic} \cap \textit{LQCoh}^{fbc}}(\mathcal{O}_\mathcal{X})
\ar[r] \ar[d]^{Rf_*} &
D^{+}_{\textit{LQCoh}^{fbc}}(\mathcal{O}_\mathcal{X})
\ar[r] \ar[d]^{Rf_*} &
D(\mathcal{O}_\mathcal{X})
\ar[d]^{Rf_*} \\
D^{+}_{\textit{Parasitic} \cap \textit{LQCoh}^{fbc}}(\mathcal{O}_\mathcal{Y})
\ar[r] &
D^{+}_{\textit{LQCoh}^{fbc}}(\mathcal{O}_\mathcal{Y}) \ar[r] &
D(\mathcal{O}_\mathcal{Y})
}
$$
을 유도하고, 따라서 몫 범주 위의 함자
$$
Rf_{\QCoh, *} :
D^{+}_\QCoh(\mathcal{O}_\mathcal{X})
\longrightarrow
D^{+}_\QCoh(\mathcal{O}_\mathcal{Y})
$$
를 유도한다. 또한 스택의 코호몰로지, 명제
\ref{stacks-cohomology-proposition-direct-image-quasi-coherent}의 함자
$R^if_\QCoh$는 (\ref{stacks-perfect-equation-Hi-quasi-coherent})의 $H^i$에 대해
$H^i \circ Rf_{\QCoh, *}$와 같다.
\end{proposition}

\begin{proof}
$E$가 $D^{+}_{\textit{LQCoh}^{fbc}}(\mathcal{O}_\mathcal{X})$에 속할 때
$Rf_*E$가
$D^{+}_{\textit{LQCoh}^{fbc}}(\mathcal{O}_\mathcal{Y})$의 대상임을
보여야 한다. 이는 스택의 코호몰로지, 명제
\ref{stacks-cohomology-proposition-loc-qcoh-flat-base-change}와 스펙트럼 열
$R^if_*H^j(E) \Rightarrow R^{i + j}f_*E$에서 따른다.
기생 가군의 경우도 스택의 코호몰로지, 보조정리
\ref{stacks-cohomology-lemma-pushforward-parasitic}를 사용하면 같은 방식으로
된다. 마지막 주장은 (\ref{stacks-perfect-equation-Hi-quasi-coherent})에 나오는
$H^i$의 정의에서 명백하다.
\end{proof}




\section{준연접 가군의 유도 당김}
\label{stacks-perfect-section-derived-pullback}

\noindent
준연접 코호몰로지 층을 갖는 복합체의 유도 당김은 일반적으로 존재한다.

\begin{proposition}
\label{stacks-perfect-proposition-derived-pullback-quasi-coherent}
$f : \mathcal{X} \to \mathcal{Y}$를 대수적 스택의 사상이라 하자.
완전 함자 $f^*$는 가환 그림
$$
\xymatrix{
D_{\textit{LQCoh}^{fbc}}(\mathcal{O}_\mathcal{X}) \ar[r] &
D(\mathcal{O}_\mathcal{X}) \\
D_{\textit{LQCoh}^{fbc}}(\mathcal{O}_\mathcal{Y})
\ar[r] \ar[u]^{f^*} &
D(\mathcal{O}_\mathcal{Y}) \ar[u]^{f^*}
}
$$
을 유도한다. 합성
$$
D_{\textit{LQCoh}^{fbc}}(\mathcal{O}_\mathcal{Y})
\xrightarrow{f^*}
D_{\textit{LQCoh}^{fbc}}(\mathcal{O}_\mathcal{X})
\xrightarrow{q_\mathcal{X}}
D_\QCoh(\mathcal{O}_\mathcal{X})
$$
은 국소화
$D_{\textit{LQCoh}^{fbc}}(\mathcal{O}_\mathcal{Y}) \to
D_\QCoh(\mathcal{O}_\mathcal{Y})$에 관해 좌유도 가능하며, 그 좌유도
함자를 $Lf^*_\QCoh$로 정의할 수 있다.
$$
Lf_\QCoh^* :
D_\QCoh(\mathcal{O}_\mathcal{Y})
\longrightarrow
D_\QCoh(\mathcal{O}_\mathcal{X})
$$
(유도 범주의 정의
\ref{derived-definition-right-derived-functor-defined}와
\ref{derived-definition-everywhere-defined} 참조.) $f$가 준콤팩트하고
준분리이면 $Lf^*_\QCoh$와 $Rf_{\QCoh, *}$는 다음 수반성을 만족한다.
$$
\Hom_{D_\QCoh(\mathcal{O}_\mathcal{X})}(Lf^*_\QCoh A, B)
=
\Hom_{D_\QCoh(\mathcal{O}_\mathcal{Y})}(A, Rf_{\QCoh, *}B)
$$
여기서 $A \in D_\QCoh(\mathcal{O}_\mathcal{Y})$이고
$B \in D^{+}_\QCoh(\mathcal{O}_\mathcal{X})$이다.
\end{proposition}

\begin{proof}
첫 번째 주장을 증명하려면
$E$가 $D_{\textit{LQCoh}^{fbc}}(\mathcal{O}_\mathcal{Y})$에 속할 때
$f^*E$가 $D_{\textit{LQCoh}^{fbc}}(\mathcal{O}_\mathcal{X})$의 대상임을
보여야 한다. $f^* = f^{-1}$가 완전하고, 스택의 코호몰로지, 명제
\ref{stacks-cohomology-proposition-loc-qcoh-flat-base-change}에 의해
$f^*$는 $\textit{LQCoh}^{fbc}(\mathcal{O}_\mathcal{Y})$를
$\textit{LQCoh}^{fbc}(\mathcal{O}_\mathcal{X})$ 안으로 보내므로
이는 즉시 따른다.

\medskip\noindent
$\mathcal{D} = D_{\textit{LQCoh}^{fbc}}(\mathcal{O}_\mathcal{Y})$로 두자.
$S$를 그 원뿔이
$D_{\textit{Parasitic} \cap \textit{LQCoh}^{fbc}}(\mathcal{O}_\mathcal{Y})$
의 대상인 $\mathcal{D}$의 사상들의 모임이라 하자.
$\mathcal{D}' = D_\QCoh(\mathcal{O}_\mathcal{X})$로 두고,
$F = q_\mathcal{X} \circ f^* : \mathcal{D} \to \mathcal{D}'$로 두자.
그러면 $\mathcal{D}, S, \mathcal{D}', F$는 유도 범주의 상황
\ref{derived-situation-derived-functor}과 정의
\ref{derived-definition-right-derived-functor-defined}에서와 같다.
$\mathcal{D}$의 임의의 대상 $E$에 대해 $LF(E)$가 정의됨을 증명하자.
즉, 보조정리 \ref{stacks-perfect-lemma-bousfield-colocalization}에서 구성한 삼각형
$$
E' \to E \to P \to E'[1]
$$
을 생각하자. $s : E' \to E$는 $S$의 원소임을 주목하라.
$E'$이 $LF$를 계산한다고 주장한다. 실제로
$s' : E'' \to E$가 $S$의 또 다른 원소라고 하자. 즉, 이는
$P'$가
$D_{\textit{Parasitic} \cap \textit{LQCoh}^{fbc}}(\mathcal{O}_\mathcal{Y})$
에 속하는 삼각형
$E'' \to E \to P' \to E''[1]$에 들어간다고 하자.
보조정리 \ref{stacks-perfect-lemma-bousfield-colocalization}(및 그 증명)에 의해
$E' \to E$는 $E'' \to E$를 거쳐 분해된다. 따라서 $E' \to E$는
계 $S/E$에서 공종말이다. 그러므로 $E'$이 $LF$를 계산함은 명백하다.

\medskip\noindent
마지막 주장을 보이기 위해 $B = q_\mathcal{X}(H)$ 및
$A = q_\mathcal{Y}(E)$로 쓰자.
위와 같이 $E' \to E$를 택한다.
한편으로는 $Rf_{\QCoh, *}(B) = q_\mathcal{Y}(Rf_*H)$를 사용하고,
다른 한편으로는 $Lf^*_\QCoh(A) = q_\mathcal{X}(f^*E')$를 사용한다.
\begin{align*}
\Hom_{D_\QCoh(\mathcal{O}_\mathcal{X})}(Lf^*_\QCoh A, B)
& = 
\Hom_{D_\QCoh(\mathcal{O}_\mathcal{X})}(q_\mathcal{X}(f^*E'),
q_\mathcal{X}(H)) \\
& = 
\colim_{H \to H'} \Hom_{D(\mathcal{O}_\mathcal{X})}(f^*E', H') \\
& = \colim_{H \to H'} \Hom_{D(\mathcal{O}_\mathcal{Y})}(E', Rf_*H') \\
& = \Hom_{D(\mathcal{O}_\mathcal{Y})}(E', Rf_*H) \\
& =
\Hom_{D_\QCoh(\mathcal{O}_\mathcal{Y})}(A, Rf_{\QCoh, *}B)
\end{align*}
여기서 쌍대극한은 그 원뿔 $P(s)$가
$D^+_{\textit{Parasitic} \cap \textit{LQCoh}^{fbc}}(\mathcal{O}_\mathcal{X})$
의 대상인
$D^+_{\textit{LQCoh}^{fbc}}(\mathcal{O}_\mathcal{X})$의 사상
$s : H \to H'$에 걸쳐 취한다.
첫 번째 등식은 위에서 보았다.
두 번째 등식은 Verdier 몫의 구성에 의해 성립한다.
세 번째 등식은 사이트 위의 코호몰로지, 보조정리
\ref{sites-cohomology-lemma-adjoint}에 의해 성립한다.
명제 \ref{stacks-perfect-proposition-derived-direct-image-quasi-coherent}에 의해
$Rf_*P(s)$는
$D^+_{\textit{Parasitic} \cap \textit{LQCoh}^{fbc}}(\mathcal{O}_\mathcal{Y})$
의 대상이므로
$\Hom_{D(\mathcal{O}_\mathcal{Y})}(E', Rf_*P(s)) = 0$이다.
따라서 네 번째 등식이 성립한다. 마지막 등식은 $E'$의 구성에 의해
성립한다.
\end{proof}







\section{유도 범주의 준연접 대상}
\label{stacks-perfect-section-QC}

\noindent
이 절은 스택 위의 층, 절 \ref{stacks-sheaves-section-QC}의 연속이다.
$\mathcal{X}$를 대수적 스택이라 하자. 그 절에서 삼각 범주
$$
\mathit{QC}(\mathcal{X}) = \mathit{QC}(\mathcal{X}_{affine}, \mathcal{O})
$$
를 정의했고, $\mathcal{X}$가 대수공간 $X$로 표현가능하면
$\mathit{QC}(\mathcal{X})$가 $D_\QCoh(\mathcal{O}_X)$와 동치임을
증명했다. 임의의 대수적 스택에 대해서도 같은 사실을 증명하기에
충분한 만큼의 이론을 정확히 마련했음이 밝혀진다.

\begin{lemma}
\label{stacks-perfect-lemma-cohomology-parasitic}
$\mathcal{X}$를 대수적 스택이라 하자. $K$를 코호몰로지 층들이
기생인 $D(\mathcal{X}_{fppf})$의 대상이라 하자. 그러면 사상
$U \to \mathcal{X}$가 평탄인 스킴 $U$ 위에 놓이는
$\mathcal{X}$의 모든 대상 $x$에 대해 $R\Gamma(x, K) = 0$이다.
\end{lemma}

\begin{proof}
절 \ref{stacks-perfect-section-lisse-etale}에서 논의한 토포스의 사상을
$g : \Sh(\mathcal{X}_{flat, fppf}) \to \Sh(\mathcal{X}_{fppf})$로 나타내자.
$x$를 사상 $U \to \mathcal{X}$가 평탄인 스킴 $U$ 위에 놓이는
$\mathcal{X}$의 대상이라 하자. 즉, $x$는
$\mathcal{X}_{flat, fppf}$의 대상이다.
보조정리 \ref{stacks-perfect-lemma-lisse-etale-cohomology}의 (2)(b)에 의해
$R\Gamma(x, K) = R\Gamma(\mathcal{X}_{flat, fppf}/x, g^{-1}K)$이다.
한편 우리의 가정은 $D(\mathcal{X}_{flat, fppf})$의 대상
$g^{-1}K$의 코호몰로지 층들이 영임을 뜻한다. 스택의 코호몰로지,
정의 \ref{stacks-cohomology-definition-parasitic}를 보라.
따라서 $g^{-1}K = 0$이고, 원하는 결론을 얻는다.
\end{proof}

\begin{lemma}
\label{stacks-perfect-lemma-cohomology-parasitic-complex}
$\mathcal{X}$를 대수적 스택이라 하자. $K$를 다음 성질을 갖는
$D(\mathcal{X}_{fppf})$의 대상이라 하자. 사상 $U \to \mathcal{X}$가
평탄인 아핀 스킴 $U$ 위에 놓이는 $\mathcal{X}$의 모든 대상 $x$에 대해
$R\Gamma(x, K) = 0$이다. 그러면 모든 $i$에 대해
$H^i(\mathcal{X}, K) = 0$이다.
\end{lemma}

\begin{proof}
절 \ref{stacks-perfect-section-lisse-etale}에서 논의한 토포스의 사상을
$g : \Sh(\mathcal{X}_{flat, fppf}) \to \Sh(\mathcal{X}_{fppf})$로 나타내자.
보조정리 \ref{stacks-perfect-lemma-lisse-etale-cohomology}의 (2)(b)에 의해, 우리의
가정은 아핀 스킴 위에 놓이는 $\mathcal{X}_{flat, fppf}$의 모든 대상
위에서 $g^{-1}K$의 코호몰로지가 소멸함을 뜻한다.
$\mathcal{X}_{flat, fppf}$의 모든 대상 $x$는 그러한 대상들에 의한
피복을 가지므로, $g^{-1}K$의 코호몰로지 층들이 소멸한다. 즉,
$g^{-1}K = 0$이다. 그러면 물론
$R\Gamma(\mathcal{X}_{flat, fppf}, g^{-1}K) = 0$이고, 이는 다시
보조정리 \ref{stacks-perfect-lemma-lisse-etale-cohomology}의 (2)(a)에 의해 원하는
결론을 준다.
\end{proof}

\begin{lemma}
\label{stacks-perfect-lemma-higher-shriek-QC}
$\mathcal{X}$를 대수적 스택이라 하자. $K$를
$D_\QCoh(\mathcal{O}_{\mathcal{X}_{flat, fppf}})$의 대상이라 하자.
그러면 $Lg_!K$는 다음 성질을 만족한다. $\mathcal{X}_{affine}$의
임의의 사상 $x \to x'$에 대해 사상
$$
R\Gamma(x', Lg_!K) \otimes_{\mathcal{O}(x')}^\mathbf{L} \mathcal{O}(x)
\longrightarrow
R\Gamma(x, Lg_!K)
$$
은 준동형이다.
\end{lemma}

\begin{proof}
보조정리 \ref{stacks-perfect-lemma-derived-quasi-coherent}의 (2)(c)에 의해 대상
$Lg_!K$는 $D_{\textit{LQCoh}^{fbc}}(\mathcal{O}_\mathcal{X})$에 속한다.
이로부터 $\mathcal{O}(x') \to \mathcal{O}(x)$가 평탄한 환 사상이면
보조정리에 표시된 사상이 동형임이 쉽게 따른다. 자세한 내용은 생략한다.

\medskip\noindent
이 문단에서는 문제가 \'etale 위상에 대해 국소적임을 논증한다.
$x \to x'$를 $\mathcal{X}_{affine}$의 일반적인 사상이라 하자.
$\{x'_i \to x'\}$를 $\mathcal{X}_{affine, \etale}$의 피복이라 하자.
$x_i = x \times_{x'} x'_i$로 두면 $\{x_i \to x\}$ 역시
$\mathcal{X}_{affine, \etale}$의 피복이다. 그러면
$\mathcal{O}(x') \to \prod \mathcal{O}(x'_i)$는 충실하게 평탄한
\'etale 환 사상이고
$$
\prod \mathcal{O}(x_i) =
\mathcal{O}(x) \otimes_{\mathcal{O}(x')} \left(\prod \mathcal{O}(x'_i)\right)
$$
이다. 따라서 우리가 생략하는 간단한 대수학 논증에 의해, 보조정리의
명제에 나오는 결과가 $\mathcal{X}_{affine}$의 각 사상
$x_i \to x'_i$에 대해 성립함을 증명하는 것으로 충분하다.
달리 말해 이 문제는 \'etale 위상에서 국소적이다.

\medskip\noindent
스킴 $X$와 전사 매끄러운 사상 $f : X \to \mathcal{X}$를 택하자.
(표기의 남용에 따라) $f$를 $\mathcal{X}$의 대상으로 볼 수 있고,
그러면 $(\Sch/X)_{fppf} = \mathcal{X}/f$이다. 스택 위의 층, 절
\ref{stacks-sheaves-section-restriction}을 보라. 예를 들어 스택 위의 층,
보조정리
\ref{stacks-sheaves-lemma-surjective-flat-locally-finite-presentation}에 의해,
$x'_i : U'_i = p(x'_i) \to \mathcal{X}$가 $f$를 거쳐 분해되게 하는
\'etale 피복 $\{x'_i \to x'\}$가 존재한다.
앞 문단의 결과에 의해, $x \to x'$가 함자
$(\Sch/X)_{fppf} \to \mathcal{X}$ 아래
$(\textit{Aff}/X)_{fppf}$의 사상 $U \to U'$의 상인 사상이라고
가정할 수 있다. 이 지점에서 $(\Sch/X)_{fppf}$로 제한한 $Lg_!K$가
보조정리 \ref{stacks-perfect-lemma-lisse-etale-functorial-derived}에 의해
$f^*Lg_!K = L(g')_!(f')^*K$와 같다는 사실을 사용함을 알 수 있다.
이로써 다음 문단에서 논의하는 경우로 귀착된다.

\medskip\noindent
$\mathcal{X} = (\Sch/X)_{fppf}$이고 $x \to x'$가 아핀 스킴의 사상
$U \to U'$에 대응한다고 하자. 여전히 $U'$ 위에서 \'etale하게
(또는 Zariski하게) 국소적으로 작업할 수 있으므로, $U' \to X$가
$X$의 어떤 아핀 열린집합을 거쳐 분해된다고 가정할 수 있다.
이로써 다음 문단에서 논의하는 경우로 귀착된다.

\medskip\noindent
$X = \Spec(R)$가 아핀 스킴이고 $\mathcal{X} = (\Sch/X)_{fppf}$이며,
$x \to x'$가 아핀 스킴의 사상 $U \to U'$에 대응한다고 하자.
$M^\bullet$을 $R\Gamma(X, K)$를 나타내는 $R$-가군의 복합체라 하자.
대수학에 관한 추가 내용, 보조정리
\ref{more-algebra-lemma-K-flat-resolution}의 구성에 의해, 각
$P_n^\bullet$이 자유 $R$-가군의 위로 유계인 복합체가 되도록
$M^\bullet = \colim P_n^\bullet$이라고 가정할 수 있다.
자세한 내용은 생략한다. 대수학에 관한 추가 내용, 비고
\ref{more-algebra-remark-P-resolution}도 보라.
규칙
$$
U \longmapsto \Gamma(U, M^\bullet \otimes_R \mathcal{O}_U)
$$
으로 주어지는 $X_{flat, fppf} = (\Sch/X)_{flat, fppf}$ 위의 가군
복합체 $M^\bullet_{flat, fppf}$를 생각하자.
강하, 절 \ref{descent-section-quasi-coherent-sheaves}의 논의에 의해
이는 층의 복합체이다. 표준 사상
$M^\bullet_{flat, fppf} \to K$가 있으며, 증명 첫머리의 언급에 의해
이 사상은 $X_{flat, fppf}$의 아핀 대상 위의 단면들에 동형을 유도한다.
$X_{flat, fppf}$의 모든 대상은 아핀 대상들에 의한 피복을 가지므로
$M^\bullet_{flat, fppf}$가 $K$와 일치함을 알 수 있다.

\medskip\noindent
$M^\bullet_{fppf}$를 위에 표시한 것과 같은 공식으로 주어지는
$X_{fppf}$ 위의 가군 복합체라 하자.
$Lg_!\mathcal{O} = g_!\mathcal{O} = \mathcal{O}$임을 상기하라.
$Lg_!$는 $g_!$의 좌유도 함자이므로
$Lg_!P_{n, flat, fppf}^\bullet = P_{n, fppf}^\bullet$을 얻는다.
함자 $Lg_!$는 호모토피 쌍대극한과 가환하고
(또는 사이트 위의 코호몰로지, 보조정리
\ref{sites-cohomology-lemma-existence-derived-lower-shriek}의 구성에 의해),
$M^\bullet = \colim P_n^\bullet$이므로
$Lg_!M^\bullet_{flat, fppf} = M^\bullet_{fppf}$을 얻는다.
$U = \Spec(A)$, $U' = \Spec(A')$이고 $U \to U'$가 환 사상
$A' \to A$에 대응한다고 하자. 위에서
$$
R\Gamma(U, Lg_!K) = M^\bullet \otimes_R A
\quad\text{및}\quad
R\Gamma(U', Lg_!K) = M^\bullet \otimes_R A'
$$
임을 알 수 있다. $M^\bullet$은 $R$-가군의 K-평탄 복합체이므로,
텐서곱의 추이성에 의해
$$
R\Gamma(U', Lg_!K) \otimes_{A'}^\mathbf{L} A
\longrightarrow
R\Gamma(U, Lg_!K)
$$
는 원하는 대로 준동형이다.
\end{proof}

\begin{proposition}
\label{stacks-perfect-proposition-QC-compare}
$\mathcal{X}$를 대수적 스택이라 하자. 그러면
$\mathit{QC}(\mathcal{X})$는 $D_\QCoh(\mathcal{O}_\mathcal{X})$와
표준적으로 동치이다.
\end{proposition}

\begin{proof}
스택 위의 층, 보조정리 \ref{stacks-sheaves-lemma-QC-compare-fppf}에 의해
비교 사상
$\epsilon : \mathcal{X}_{affine, fppf} \to \mathcal{X}_{affine}$에
의한 당김은 $\mathit{QC}(\mathcal{X})$를 충만한 부분범주
$Q_\mathcal{X} \subset D(\mathcal{X}_{affine, fppf}, \mathcal{O})$와
동일시한다. 스택 위의 층, 식
(\ref{stacks-sheaves-equation-alternative-ringed})에 나오는
환 달린 토포스의 동치를 사용하여, $Q_\mathcal{X}$를
$D(\mathcal{X}_{fppf}, \mathcal{O})$의 충만한 부분범주로 볼 수 있고
그렇게 보기로 한다.

\medskip\noindent
마찬가지로 보조정리 \ref{stacks-perfect-lemma-bousfield-colocalization}와 비고
\ref{stacks-perfect-remark-QCoh-admissible}에 의해, $D_\QCoh(\mathcal{O}_\mathcal{X})$를
$D_{\textit{LQCoh}^{fbc}}(\mathcal{O}_\mathcal{X})$의 왼쪽 허용 가능
부분범주
$D_{\textit{Parasitic} \cap \textit{LQCoh}^{fbc}}(\mathcal{O}_\mathcal{X})$
의 왼쪽 직교 $\mathcal{A}$로 볼 수 있다.

\medskip\noindent
증명을 끝내기 위해 $D(\mathcal{X}_{fppf}, \mathcal{O})$의 부분범주로서
$Q_\mathcal{X}$와 $\mathcal{A}$가 같음을 보이겠다.

\medskip\noindent
1단계: $Q_\mathcal{X}$는
$D_{\textit{LQCoh}^{fbc}}(\mathcal{O}_\mathcal{X})$에 포함된다.
$Q_\mathcal{X}$의 대상 $K$는 다음 성질로 특징지어진다.
$K$를 $D(\mathcal{X}_{affine, fppf}, \mathcal{O})$의 대상으로 볼 때,
$R\epsilon_*K$는 $\mathit{QC}(\mathcal{X}_{affine}, \mathcal{O})$의
대상이다. 이는 다시 정확히 $\mathcal{X}_{affine}$의 모든 사상
$x \to x'$에 대해 사상
$$
R\Gamma(x', K) \otimes_{\mathcal{O}(x')}^\mathbf{L} \mathcal{O}(x)
\longrightarrow
R\Gamma(x, K)
$$
이 동형임을 뜻한다. 사이트 위의 코호몰로지, 보조정리
\ref{sites-cohomology-lemma-cartesion-plus-topology}의 명제에 붙은
각주를 보라. 이제 $x' \to x$가 아핀 스킴의 평탄 사상 위에 놓이면,
이는
$$
H^i(x', K) \otimes_{\mathcal{O}(x')} \mathcal{O}(x)
\cong
H^i(x, K)
$$
임을 뜻한다. 이는 분명히 $H^i(K)$가 \'etale 위상에 대한 층이고
(스택 위의 층, 보조정리
\ref{stacks-sheaves-lemma-quasi-coherent-alternative}), 평탄 밑변환
성질을 가짐을 뜻한다(작은 세부 사항은 생략한다).

\medskip\noindent
2단계: $Q_\mathcal{X}$는 $\mathcal{A}$에 포함된다.
이를 보이려면 $K$가 $Q_\mathcal{X}$에 속할 때, 모든 $P$가
$D_{\textit{Parasitic} \cap \textit{LQCoh}^{fbc}}(\mathcal{O}_\mathcal{X})$
에 속하는 경우에 대해 $\Hom(K, P) = 0$임을 보이는 것으로
충분하다. 대상
$$
H = R\SheafHom_{\mathcal{O}_\mathcal{X}}(K, P)
$$
를 생각하자. $x$를 아핀 스킴 $U = p(x)$ 위에 놓이는
$\mathcal{X}$의 대상이라 하자. 사이트 위의 코호몰로지, 보조정리
\ref{sites-cohomology-lemma-section-RHom-over-U}에 의해 다음 식의
첫 번째 등식을 얻는다.
$$
R\Gamma(x, H) =
R\Hom_{\mathcal{O}_\mathcal{X}}(K|_{\mathcal{X}/x}, P|_{\mathcal{X}/x}) =
R\Hom_{\mathcal{O}}(K|_{\mathcal{X}_{affine}/x}, P|_{\mathcal{X}_{affine}/x})
$$
두 번째 등식은 사이트 $\mathcal{X}/x$의 토포스가 사이트
$\mathcal{X}_{affine}/x$의 토포스와 동치라는 사실에서 나온다.
스택 위의 층, 식
(\ref{stacks-sheaves-equation-alternative-ringed})을 보라.
어떤 $N$이 $\mathit{QC}(\mathcal{O})$에 속하여
$K = \epsilon^*N$으로 쓸 수 있다. 그러면 사이트 위의 코호몰로지,
보조정리 \ref{sites-cohomology-lemma-QC-hom-out-of-plus-topology}에 의해
$$
R\Gamma(x, H) =
R\Hom_{D(\mathcal{O}(x))}(R\Gamma(x, N), R\Gamma(x, P))
$$
임을 알 수 있다. 보조정리 \ref{stacks-perfect-lemma-cohomology-parasitic}에 의해
$U \to \mathcal{X}$가 평탄하면 $R\Gamma(x, P) = 0$이고, 따라서 같은
가정 아래 $R\Gamma(x, H) = 0$이다. 보조정리
\ref{stacks-perfect-lemma-cohomology-parasitic-complex}에 의해
$R\Gamma(\mathcal{X}, H) = 0$이고, 그러므로 $\Hom(K, P) = 0$이다.

\medskip\noindent
3단계: $\mathcal{A}$는 $Q_\mathcal{X}$에 포함된다.
$K$를 $\mathcal{A}$의 대상이라 하고 $x \to x'$를
$\mathcal{X}_{affine}$의 사상이라 하자. 사상
$$
R\Gamma(x', K) \otimes_{\mathcal{O}(x')}^\mathbf{L} \mathcal{O}(x)
\longrightarrow
R\Gamma(x, K)
$$
이 준동형임을 보여야 한다. 사이트 위의 코호몰로지, 보조정리
\ref{sites-cohomology-lemma-cartesion-plus-topology}의 명제에 붙은 각주를
보라. 보조정리 \ref{stacks-perfect-lemma-bousfield-colocalization}의 증명과 비고
\ref{stacks-perfect-remark-QCoh-admissible}의 논의에 의해, $\mathcal{A}$는
$D_\QCoh(\mathcal{O}_{\mathcal{X}_{flat, fppf}})$로 제한한 $Lg_!$의 상이다.
따라서 어떤 $M$이 $D_\QCoh(\mathcal{O}_{\mathcal{X}_{flat, fppf}})$에
속하여 $K = Lg_!M$이라고 가정할 수 있다.
그러면 원하는 등식은 보조정리 \ref{stacks-perfect-lemma-higher-shriek-QC}에서 따른다.
\end{proof}
